\documentclass[12pt]{article}

\usepackage[utf8]{inputenc}
\usepackage{amsmath, amssymb, amsthm}
\usepackage{geometry}
\geometry{a4paper, margin=2.5cm}

\title{Teorema lui Goursat}
\author{}
\date{}

\theoremstyle{definition}
\newtheorem{definition}{Definiție}

\theoremstyle{plain}
\newtheorem{theorem}{Teoremă}
\newtheorem{proposition}{Propoziție}

\begin{document}

\maketitle

\begin{definition}
O funcție $f:\Omega \to \mathbb{C}$ se numește \emph{olomorfă} într-un domeniu $\Omega$ dacă este derivabilă complex în fiecare punct al lui $\Omega$.
\end{definition}

\begin{definition}
O \emph{curbă închisă, netedă, orientată pozitiv} este o parametrizare


\[
\gamma : [a,b] \to \Omega
\]


care este continuă pe $[a,b]$, derivabilă pe $(a,b)$, satisface $\gamma(a)=\gamma(b)$ și are $\gamma'(t)\neq 0$ pentru toate $t\in(a,b)$.
\end{definition}

\begin{theorem}
(\textbf{Teorema lui Goursat})  
Fie $\Omega \subset \mathbb{C}$ un domeniu și fie $f:\Omega \to \mathbb{C}$ o funcție olomorfă.  
Atunci, pentru orice triunghi $T$ complet conținut în $\Omega$, are loc


\[
\int_{\partial T} f(z)\,dz = 0.
\]


Cu alte cuvinte, integrala unei funcții olomorfe pe conturul unui triunghi este zero, \emph{fără a presupune continuitatea derivatei} lui $f$.
\end{theorem}

\begin{proposition}
(\textbf{Consecință})  
Dacă $f$ este olomorfă într-un domeniu $\Omega$, atunci pentru orice curbă închisă netedă $\gamma$ cu imaginea conținută în $\Omega$,


\[
\int_{\gamma} f(z)\,dz = 0.
\]


Aceasta este forma generalizată a teoremei lui Cauchy, obținută ca urmare a teoremei lui Goursat.
\end{proposition}

\end{document}
